% Part: turing-machines
% Chapter: machines-computations
% Section: turing-machines

\documentclass[../../../include/open-logic-section]{subfiles}

\begin{document}

\olfileid{tur}{mac}{tur}
\olsection{ట్యూరింగ్ యంత్రాలు}

\begin{explain}
దేనిని ట్యూరింగ్ యంత్రం అనాలో చెప్పే ఆచారబద్ధ నిర్వచనం అమూర్తంగా
కనిపిస్తుంది; కానీ నిజానికి అది సరళమైనదే: ట్యూరింగ్ యంత్రపు పనితీరును
నిర్దేశించడానికి అవసరమైన సమాచారమంతటినీ ఒకే గణిత నిర్మాణంలో అది కూర్చుతుంది.
ఈ సమాచారంలో (1) యంత్రం ఏ స్థితుల్లో ఉండగలదు, (2) టేపుపై ఏ సంకేతాలకు అనుమతి
ఉంది, (3) యంత్రం ఏ స్థితిలో ప్రారంభం కావాలి, (4) యంత్రపు నిర్దేశ సమితి ఏమిటి
అనేవి ఉంటాయి.
\end{explain}

\begin{defn}[ట్యూరింగ్ యంత్రం]
\emph{ట్యూరింగ్ యంత్రం} $M$ అనేది కింది వాటితో కూడిన
$\langle Q, \Sigma, q_0, \delta\rangle$ అనే జతక్రమం:
\begin{enumerate}
\item పరిమిత \emph{స్థితుల} సమితి~$Q$,
\item $\TMendtape$, $\TMblank$లను కలిగి ఉండే పరిమిత \emph{వర్ణమాల} $\Sigma$,
\item ఒక \emph{ప్రారంభ స్థితి}~$q_0 \in Q$,
\item పరిమిత \emph{నిర్దేశ సమితి}~$\delta\colon Q \times \Sigma
  \pto Q \times \Sigma \times \{\TMleft, \TMright, \TMstay\}$.
\end{enumerate}
పాక్షిక ప్రమేయం~$\delta$ను $M$ యొక్క \emph{సంక్రమణ ప్రమేయం} అని కూడా అంటారు.
\end{defn}

\begin{explain}
టేపు ఒకే దిశలో అనంతంగా ఉంటుందని ఊహిస్తాం. అందువల్ల టేపు యొక్క ఎడమ చివర
గుర్తుగా ప్రత్యేక సంకేతం~$\TMendtape$ను నిర్దేశించడం ఉపయోగకరం. ట్యూరింగ్ యంత్ర
ప్రోగ్రాములు టేపు నుండి ``బయటకు వెళ్లే ప్రమాదంలో'' ఎప్పుడు ఉన్నాయో తెలుసుకోవడం
దీనితో సులభమవుతుంది. ఈ సంకేతంపై మరొకటి ఎన్నటికీ వ్రాయరాదని, అంటే
$\delta(q,\TMendtape)$ నిర్వచితమైతే
$\delta(q,\TMendtape) = \tuple{q', \TMendtape, D}$ అని ఊహించవచ్చు; ఇక్కడ
$D \in \{\TMleft,\TMright,\TMstay\}$. కొన్ని
పాఠ్యపుస్తకాలు ఇలా చేస్తాయి; మనం చేయం. మీ ట్యూరింగ్ యంత్రాన్ని నిర్మించేటప్పుడు
అది $\TMendtape$పై ఎన్నటికీ మరొకటి వ్రాయకుండా జాగ్రత్తపడవచ్చు. అంతేకాకుండా,
అలా వ్రాయడానికి అనుమతించడం ఉపయోగకరమైన వెసులుబాటును ఇచ్చే సందర్భాలూ ఉన్నాయి.
\sourcecorrection{OLTETURMAC-002}{ఆంగ్ల మూలం టేపు-అంత సంకేతాన్ని కాపాడే త్రయంలో దిశను నిర్వచించని $x$తో సూచించింది. ఈ అధ్యాయం అంతటా దిశకు వాడిన $D$ను పునరుద్ధరించి, దాని మూడు సాధ్యమైన విలువలను స్పష్టం చేశాం.}
\end{explain}

\begin{ex}
\emph{సరి యంత్రం:} ఆచారబద్ధంగా సరి యంత్రం
$\tuple{Q, \Sigma, q_0, \delta}$ అనే చతుష్టయం; ఇక్కడ
\begin{align*}
Q & = \{ q_0, q_1 \} \\
\Sigma & = \{ \TMendtape, \TMblank, \TMstroke \}, \\
\delta(q_0, \TMstroke) & = \tuple{q_1, \TMstroke, \TMright},\\
\delta(q_1, \TMstroke) & = \tuple{q_0, \TMstroke, \TMright},\\
\delta(q_1, \TMblank)  & = \tuple{q_1, \TMblank, \TMright}.
\end{align*}
\end{ex}

\end{document}
